\SourceLeaf{68}
Damit ist die gestellte Aufgabe gelöst, soweit dies durch Reihen möglich war.

Der hochangesehene \textsc{Kl\"ugel} hat eine verwandte Untersuchung vorgelegt (über die Potenzialfunktion

\[
(1-2x\cos\varphi+xx)^{-m}
\]

\noindent in Commentat. Societatis huius Tom. XII). Die von ihm betrachtete Form
weicht zwar von derjenigen ab, die in dieser kleinen Abhandlung nach \textsc{Euler}s Vorbild untersucht wurde.
Es wäre jedoch leicht, die eine aus der anderen herzuleiten. Damit es nicht so aussieht,
als hätte ich bloß bereits Geleistetes wiederholt, möchte ich darauf hinweisen, dass ich die ganze Untersuchung
unternommen habe, um an einem Beispiel zu zeigen, dass die Behandlung unendlicher Reihen durch ihre allgemeinen Glieder
keine Schwierigkeiten bereitet und oft leichter als besondere Kunstgriffe zur Auffindung allgemeiner Gesetze führt.
Im Hauptteil dieser kleinen Abhandlung habe ich daher einen anderen Weg eingeschlagen
als die hochangesehenen \textsc{Euler} und \textsc{Kl\"ugel}. Im Übrigen habe ich versucht, alles auf allgemeine Ausdrücke zurückzuführen
und, wo immer möglich, mit einem Beweis zu versehen.

\textls[140]{Korollar.}

11. In der höheren Analysis kommen bekanntlich endliche wie unendliche Reihen vor, deren Glieder
auf den ersten Blick verschiedenen Gesetzen zu folgen scheinen. Das gilt insbesondere für Reihen, in
denen die Glieder mit geradem Index anders gebildet sind als die mit ungeradem.
Üblicherweise umschreibt man dann jede der beiden Arten durch ein eigenes Gesetz. Außer ihrer Weitschweifigkeit
haben derartige Definitionen etwas an sich, das der Natur und Leistungsfähigkeit der Analysis zuwiderläuft. Ihr Vorzug
besteht gerade darin, selbst höchst unterschiedliche Werte einer Größe, wenn solche vorkommen,
in einem einzigen allgemeinen Ausdruck zusammenzufassen und sie alle daraus herzuleiten.
Wenn also bei der Lösung einer geometrischen Aufgabe mehrere Werte ein und derselben Größe eine Lösung liefern,
so sind sie alle in der Gleichung enthalten, durch die wir die Beziehung der gegebenen Größen
zur gesuchten ausdrücken. Etwas Ähnliches lässt sich von den Ausdrücken verlangen, die die allgemeinen
Glieder von Reihen darstellen. Unterscheiden sich die geraden und
ungeraden Glieder, so muss die Formel für das allgemeine Glied, wenn sie vollständig sein soll,
für einen unbestimmten geraden Gliedindex die einen und für einen
ungeraden die anderen Werte liefern. Wie dies erreicht werden kann, habe ich bereits in der Einleitung
zu dieser kleinen Abhandlung dargelegt. Für $(-1)^m$ gilt nämlich bei geradem $m$ der Wert $=+1$, bei ungeradem der Wert $=-1$. Somit
